\documentclass[10pt]{article}

\usepackage[letterpaper,margin=0.42in]{geometry}
\usepackage{amsmath,amssymb}
\usepackage{multicol}

\pagestyle{empty}
\setlength{\parindent}{0pt}
\setlength{\parskip}{1.4pt}
\setlength{\columnsep}{0.25in}
\setlength{\abovedisplayskip}{2pt}
\setlength{\belowdisplayskip}{2pt}

\newcommand{\cube}[1]{\{0,1\}^{#1}}
\newcommand{\NAE}{\operatorname{NAE}}
\newcommand{\ANF}{\operatorname{ANF}}

\begin{document}

\begin{center}
  {\large\bfseries Theorem Summary: Exact Repair Obstructions for a\\[-1pt]
  Sensitivity--14 Kushilevitz Composition}\par
  \vspace{2pt}
  {\small Jia Lei and Bokun Zhao \quad August 14, 2026}
\end{center}

\vspace{-5pt}
\scriptsize
\begin{multicols}{2}

\textbf{Setup and historical base.}
Let
\[
 K(y)=\sum_i y_i-\sum_{i<j}y_iy_j+\sum_{T\in\mathcal D}\prod_{i\in T}y_i,
\]
\[
\mathcal D=\{013,014,024,025,035,345,145,125,123,234\}.
\]
The set $\mathcal D$ is the $2\!-\!(6,3,2)$ design, and $K$ is a relabeling
of the degree-three, sensitivity-six function of Kushilevitz recorded by
Nisan and Wigderson.  For four disjoint raw triples $X_i$, put
\[
 L_i=\sum_{a\in X_i}x_a,\qquad
 E_i=\sum_{\{a,b\}\subset X_i}x_ax_b,\qquad
 g_i=\NAE_3(X_i)=L_i-E_i.
\]

\textbf{Near-solution (degree \(6\), sensitivity \(14\)).}
The Boolean function
\[
 P=K(g_0,g_1,g_2,g_3,x_{12},x_{13})
\]
satisfies $P(0)=0$ and $P(e_j)=1$ for $0\le j<14$.  Thus
$s(P,0)=14$.  Its complete highest layer is
\[
 P_6=-E_1E_3(E_0+E_2).
\]
It has exactly $2\cdot3^3=54$ distinct degree-six monomials, all with
coefficient $-1$, and none of higher degree.

\textbf{Two noncancellation principles.}
First uniquely multilinearize a nonzero meta-polynomial $Q$.  Under substitution of
disjoint NAE blocks and disjoint coordinate ports, distinct meta monomials
have distinct raw top signatures.  Giving NAE variables weight $2$ and ports
weight $1$,
\[
 \deg Q(g,z)=\max_{c_{I,J}\ne0}\bigl(2|I|+|J|\bigr).
\]
Disjointness and multilinearization are essential.

For the separate Walsh lemma, in spin variables let
\[
 \chi_i=s_{i0}s_{i1}s_{i2},\quad
 u_i=\frac{s_{i0}s_{i1}+s_{i0}s_{i2}+s_{i1}s_{i2}}2,\quad
 r_i=1-2g_i=u_i-\frac12.
\]
If $q=\sum_{U,B}d_{U,B}u_Ut_B$, multiplication by
$\chi_0\chi_1\chi_2\chi_3$ sends each nonzero sector $(U,B)$ to disjoint
Walsh characters of degree
\[
 12-2|U|+|B|.
\]
This block-signature statement is not an application of the weighted lemma.

\textbf{Central selector obstruction.}
Define
\[
 A=K(g_0,g_1,g_2,x_9,x_{10},x_{11}).
\]
Then $A$ is Boolean, has degree $5$, sensitivity $12$, and
\[
 A_5=E_0E_1(x_9+x_{10})+E_0E_2(x_{10}+x_{11})
      +E_1E_2(x_{11}+x_9).
\]
There is no Boolean $B:\cube{12}\to\{0,1\}$ with
$B(0)=1$ and $\deg(B-A)\le4$.  Averaging $B$ over block permutations and
whole-block complements gives a $[0,1]$-valued function of the NAE outputs and
ports.  Its fixed-port heavy slices are quadratic.  The sums $S_z$ of their
three pair coefficients satisfy
\[
 S_{111}=S_{000}+6,\quad
 S_{000}\ge-3+2C(0),\quad S_{111}\le3,\quad C(0)\ge\tfrac18,
\]
a contradiction.  Here $B$ may depend arbitrarily on all twelve raw variables;
the theorem is scoped to the fixed base $A$.

\textbf{Gate-factor and fixed-routing obstructions.}
Let $T:\cube6\to\{0,1\}$ and
$F=T(g_0,g_1,g_2,g_3,x_{12},x_{13})$.  If $F(0)=0$ and
$F(e_j)=1$ for every $0\le j<14$, then $\deg F>5$: setting the ports to zero
would leave a Boolean quadratic with four sensitive meta-bits, whereas every
Boolean quadratic has sensitivity at most $3$.  Moreover, let each $h_i$ be a
Boolean function of degree at most two with $h_i(0)=0$, sensitive at its own
origin in all three coordinates of a designated pairwise-disjoint triple.
Routing $h_i$ to input $i$ of $K$, with two distinct coordinate ports disjoint
from all triples routed to inputs $4,5$, cannot evade the degree-six layer.

\textbf{Full-parity-twist obstruction.}
For every Boolean meta-function $T$,
\[
 T(g_0,g_1,g_2,g_3,x_{12},x_{13})
 \oplus p_0\oplus p_1\oplus p_2\oplus p_3
\]
has degree $>5$, where $p_i$ is raw XOR parity on block $i$.  The Walsh lemma
leaves only a putative output spin
$q_T=a(t_{12},t_{13})\prod_i(r_i+1/2)$.  If $a=0$, Booleanity fails; if
$a\ne0$, its magnitude changes by factors of $3$ between all-equal and NAE
block inputs, contradicting $|q_T|=1$.  Partial twists and raw
orientation-dependent $T$ are outside the theorem.

\textbf{Exact additive reformulation.}
Let $\mathcal T=\{0,e_0,\ldots,e_{13}\}$ and
$\Delta=E_1E_3(E_0+E_2)=-P_6$.  For real multilinear $R$, $F=P+R$ is
Boolean, has degree at most $5$, and agrees with $P$ on $\mathcal T$ if and
only if
\[
 R=\Delta+L\ (\deg L\le5),\quad
 R(R+2P-1)=0\text{ on }\cube{14},\quad R|_{\mathcal T}=0.
\]
Every admissible $F$ occurs exactly once, with $R=F-P$.

\textbf{Exact XOR reformulation.}
Every Boolean $F$ occurs uniquely as
\[
 J=P\oplus F,\qquad F=P+(1-2P)J.
\]
The desired functions correspond exactly to Boolean masks $J$ with
$J|_{\mathcal T}=0$ and $\deg_{\mathbb R}(P+(1-2P)J)\le5$.  Necessarily, but
not sufficiently,
$\ANF_{\ge6}(J)=\ANF_{\ge6}(P)$.  The additive and XOR systems are lossless
changes of coordinates, not restricted search families.

\textbf{The simplest mask fails at degree \(10\).}
For $J_0=g_1g_3(g_0\oplus g_2)$,
\[
 \deg_{\mathbb R}J_0=8,\quad (J_0)_8=-2E_0E_1E_2E_3,
\]
whereas exact integer M\"obius inversion gives
\[
 \deg(P\oplus J_0)=10,\quad
 (P\oplus J_0)_{10}=4E_0E_1E_2E_3x_{12}x_{13}.
\]
The top defect has exactly $3^4=81$ monomials, all coefficient $+4$.

\textbf{Status and review focus.}
No degree-five, sensitivity-fourteen construction, improved recursive
exponent, or priority claim is made.  The accompanying standard-library
verifier reconstructs $K,A,P,J_0,P\oplus J_0$ and checks their exact integer
M\"obius coefficients.  Specialist review should focus on both signature
lemmas, the selector averaging, gate-factor and parity scopes, both exact
systems, the degree-ten calculation, and the literature comparison.

\end{multicols}
\end{document}
